% 355_SSB_Torus_En_ch.tex
% Johann Pascher, 7 September 2026
\providecommand{\BM}{\textbf{[B]}}
\providecommand{\KM}{\textbf{[K]}}
\providecommand{\SM}{\textbf{[S]}}
\providecommand{\XM}{\textbf{[X]}}
\providecommand{\QM}{\textbf{[Q]}}

\phantomsection
\addcontentsline{toc}{chapter}{%
  SSB $SU(2)_L\times U(1)_Y\to U(1)_\text{EM}$ from $T^4/Z_3$ Geometry}

\section*{Status Markers}
\begin{center}
\begin{tabular}{@{}lp{10.5cm}@{}}
\textbf{[K]} & \emph{Confirmed} --- numerically verified. \\[3pt]
\textbf{[B]} & \emph{Proven} --- algebraically demonstrated. \\[3pt]
\textbf{[S]} & \emph{Speculative} --- open. \\[3pt]
\textbf{[X]} & \emph{Experimental} --- confirmed. \\[3pt]
\textbf{[Q]} & \emph{Source} --- external primary value. \\
\end{tabular}
\end{center}

\section*{Abstract}

The spontaneous symmetry breaking $SU(2)_L\times U(1)_Y\to U(1)_\text{EM}$
is implemented in the Standard Model through the Higgs potential $V(\Phi)$
with $\mu^2<0$ --- a postulate. This document derives it from three
building blocks of the $T^4/Z_3$ geometry, all already proven in the corpus.

\textbf{The three building blocks:}
\begin{enumerate}
\item $\tilde{T}\cdot m=1$ forces $\langle\Phi\rangle\neq0$: the vacuum
  must carry mass \BM\ (Docs.~003, 010, 306, 334).
\item The $Z_3$ Galois structure forces $Q_\text{vac}=0$: the vacuum
  is electrically neutral \BM\ (Doc.~346).
\item $Q=0$ of the vacuum selects $U(1)_\text{EM}$ as the residual group:
  the Gell-Mann-Nishijima generator annihilates the vacuum \BM\ (Doc.~347).
\end{enumerate}

\textbf{Main finding:} The symmetry-breaking structure is topologically forced,
not postulated \BM. What appears in SM convention as the spontaneously breaking
postulate $\mu^2<0$ is in FFGFT a consequence of the time-field geometry and
the Galois classification.

\textbf{Open point \SM:} CKM/PMNS mixing angle values (Doc.~348 Theorem~D).
New candidate: GF$(3^6)$ = GF$(729)$ and 7th roots of unity
(Matzke IPI mail, 7~September 2026, R115).

\textbf{Overview:} All SM fermion quantum numbers and their certificate
status are summarised in Doc.~356 (particle diagram) and Doc.~357
(certificate table).

\section{Background}

\subsection*{The SM Higgs Potential}

The Standard Model implements SSB through:
\begin{equation}
  V(\Phi) = \mu^2\,|\Phi|^2 + \lambda\,|\Phi|^4,
  \quad \mu^2 < 0,
  \label{eq:higgs_pot}
\end{equation}
with minimum $|\Phi|^2 = -\mu^2/(2\lambda) = v^2/2$, $v\approx246\,\text{GeV}$.
The Higgs doublet $\Phi = (\phi^+,\phi^0)^T$ has $I_3 = (-\tfrac{1}{2},+\tfrac{1}{2})$,
$Y = +1$. The neutral component develops the VEV:
\begin{equation}
  \langle\Phi\rangle = \begin{pmatrix}0\\v/\sqrt{2}\end{pmatrix},
  \quad Q_\text{vac} = I_3 + \frac{Y}{2} = -\tfrac{1}{2}+\tfrac{1}{2} = 0.
  \label{eq:sm_vev}
\end{equation}
The surviving symmetry is $U(1)_Q$ with generator $Q=I_3+Y/2$ ---
this is $U(1)_\text{EM}$ \XM\ (PDG~\QM).

\textbf{The question:} Why does the vacuum have $Q=0$? Why $\mu^2<0$?
In the SM: postulated. In FFGFT: from the geometry.

\section{Building Block 1: VEV $\neq 0$ from $\tilde{T}\cdot m = 1$}

The time-field duality (Docs.~003, 010) in natural units:
\begin{equation}
  \tilde{T}(x,t)\cdot m(x,t) = 1.
  \label{eq:tm1}
\end{equation}
This relation has no solution at $m=0$: $\tilde{T}\to\infty$ is not a
physical state (Docs.~306, 334). The vacuum must carry a non-vanishing mass:
\begin{equation}
  m_\text{vac} \neq 0 \quad\Longleftrightarrow\quad \langle\Phi\rangle \neq 0.
  \label{eq:vev_nonzero}
\end{equation}

\textbf{Connection to the Higgs potential:} In SM language, $m_\text{vac}=0$
is the minimum of $V$ for $\mu^2>0$ --- an empty vacuum. For $\mu^2<0$
the minimum lies at $|\Phi|\neq0$. The condition \eqref{eq:vev_nonzero}
from FFGFT is equivalent to the requirement $\mu^2<0$ in the SM potential:
both force $\langle\Phi\rangle\neq0$ \BM.

\textbf{The difference:} In the SM, $\mu^2<0$ is a postulate.
In FFGFT, $m_\text{vac}\neq0$ follows from the reciprocity
$\tilde{T}\cdot m=1$, which itself follows from $E=mc^2$
(Doc.~354 Insight~1) \BM.

\section{Building Block 2: $Q_\text{vac} = 0$ from the Galois Structure}

\subsection*{Step 2a: The vacuum is a torus background}

The time field in the vacuum is $\tilde{T}_0 = 1/m_0 = \text{const}$ --- a
static, spatially homogeneous configuration. In the $T^4/Z_3$ framework
this is the $k=0$ mode (zero winding): no angular momentum, no colour charge,
no electric charge. The vacuum is topologically trivial \BM\ (Doc.~A075).

\subsection*{Step 2b: Galois classification of the vacuum}

In GF$(27)^*$, the combination of QR-property mod~13 and $N\bmod3$
colour charge classifies all fermions (Doc.~346 Theorem~B \BM):

\begin{center}
\begin{tabular}{@{}llll@{}}
\toprule
QR mod~13 & $N\bmod3$ & Particle type & $Q$ \\
\midrule
QR (Quadratic Residue) & $\equiv 0$ & Neutrino & $0$ \\
NQR (Non-QR) & $\equiv 0$ & Charged lepton & $-1$ \\
NQR & $\equiv 1,2$ & Quark & $\pm1/3,\pm2/3$ \\
\bottomrule
\end{tabular}
\end{center}

The vacuum is no fermion --- it is the $k=0$ mode, the topologically trivial
element. Its decisive property:
\begin{itemize}
  \item No colour charge: $N\bmod3 = 0$ (colour singlet) \BM
  \item Electrically neutral: QR-property, hence $Q=0$ \BM
\end{itemize}

This follows from Doc.~346 Theorem~A \BM:
$Q=0 \iff \text{QR mod~13}$.
The topologically trivial element (zero winding) is QR \BM.

\subsection*{Step 2c: The vacuum has $I_3=-1/2$, $Y=+1$ with $Q=0$}

The doublet structure follows from the Sd sector (Doc.~349 Theorem~D \BM):
the Sd sector (left-handed, odd Cl$(6)$ grades) carries $I_3=-1/2$.
The hypercharge $Y$ follows from the Legendre symbol (Doc.~347 Theorem~A \BM).
The vacuum in the Sd sector with $I_3=-1/2$, $Y=+1$:
\begin{equation}
  Q_\text{vac} = I_3 + \frac{Y}{2} = -\frac{1}{2} + \frac{1}{2} = 0.
  \label{eq:qvac}
\end{equation}
$Q_\text{vac}=0$ is algebraically forced \BM.

\section{Building Block 3: $U(1)_\text{EM}$ as Residual Group}

\subsection*{Residual symmetry of the vacuum}

The symmetry group $G = SU(2)_L\times U(1)_Y$ acts on the vacuum.
The residual group $H\subset G$ after breaking is the subgroup that
leaves the vacuum invariant:
\begin{equation}
  H = \{g\in G \mid g\,|\text{vac}\rangle = |\text{vac}\rangle\}.
\end{equation}

The generator $Q = I_3 + Y/2$ (Gell-Mann-Nishijima, Doc.~347 Theorem~C \BM)
annihilates the vacuum:
\begin{equation}
  Q\,|\text{vac}\rangle = Q_\text{vac}\,|\text{vac}\rangle = 0.
\end{equation}
The one-parameter subgroup:
\begin{equation}
  U(1)_Q : \quad e^{i\alpha Q}\,|\text{vac}\rangle = |\text{vac}\rangle
  \quad\forall\,\alpha\in\mathbb{R}.
\end{equation}
$U(1)_Q$ leaves the vacuum invariant --- this is $U(1)_\text{EM}$ \BM.

\subsection*{All other generators break the symmetry}

The remaining generators of $SU(2)_L\times U(1)_Y$ generate the massive
bosons $W^\pm$, $Z^0$ (Goldstone theorem \XM):
\begin{align}
  W^\pm &\leftrightarrow T_\pm = T_1 \pm iT_2, \quad T_\pm|\text{vac}\rangle \neq 0 \\
  Z^0 &\leftrightarrow T_3 - \frac{Y}{2}\sin^2\theta_W, \quad Z^0|\text{vac}\rangle \neq 0.
\end{align}

The photon $\gamma$ corresponds to the generator $Q = T_3 + Y/2$, which
annihilates the vacuum --- it remains massless \BM.
The mass ratio $M_W/M_Z = \cos\theta_W$ follows from the Weinberg angle
derived from $\xi$ (Doc.~323 \KM).

\section{Complete Chain}

\begin{center}
\begin{tabular}{@{}p{0.6cm}p{7.5cm}p{1.5cm}p{3.5cm}@{}}
\toprule
Step & Statement & Status & Source \\
\midrule
1 & $\tilde{T}\cdot m=1$ from $E=mc^2$ & \BM & Docs.~003,010,052 \\
2 & $m_\text{vac}\neq0$: no massless vacuum & \BM & Docs.~306,334 \\
3 & Vacuum is $k=0$ mode (topologically trivial) & \BM & Doc.~A075 \\
4 & $Q_\text{vac}=0$: trivial element is QR & \BM & Doc.~346 Thm.~A \\
5 & Sd sector: $I_3=-1/2$ (left-handed) & \BM & Doc.~349 Thm.~D \\
6 & $Y=+1$ from Legendre symbol & \BM & Doc.~347 Thm.~A \\
7 & $Q=I_3+Y/2=0$ (Gell-Mann-Nishijima) & \BM & Doc.~347 Thm.~C \\
8 & $U(1)_Q$ leaves vacuum invariant: residual group & \BM & this doc. \\
9 & $W^\pm$, $Z^0$ massive; $\gamma$ massless & \BM+\XM & Goldstone + Doc.~323 \\
\midrule
\multicolumn{2}{@{}l}{SSB $SU(2)_L\times U(1)_Y\to U(1)_\text{EM}$} & \BM & complete \\
\bottomrule
\end{tabular}
\end{center}

\section{Numerical Consistency}

\begin{center}
\begin{tabular}{@{}lllll@{}}
\toprule
Quantity & FFGFT & Experiment \QM & Dev. & Status \\
\midrule
$\sin^2\theta_W(M_Z)$ & $0.2308$ & $0.2312$ & $-0.17\,\%$ & \KM \\
$M_W/M_Z$ & $\cos\theta_W = 0.8770$ & $0.8814$ & $-0.50\,\%$ & \KM \\
$Q_\text{vac}$ & $0$ (exact) & $0$ (exact) & $0$ & \BM \\
$m_\gamma$ & $0$ (exact) & $<10^{-18}$\,eV & --- & \BM+\XM \\
\bottomrule
\end{tabular}
\end{center}

\section{Open Points}

\textbf{CKM/PMNS mixing angle values \SM:}
The structure of the mixing matrices (which entries vanish, which are
$1/\sqrt{2}$) is proven \BM\ (Doc.~348 Theorems B,C). The numerical
values of the mixing angles require further algebraic work
(Doc.~348 Theorem~D). This is the only remaining \SM\ in the
gauge/Higgs complex.

\textbf{New candidate: GF$(3^6)$ = GF$(729)$ and 7th roots of unity \SM:}

IPI correspondence with Douglas Matzke (7~September 2026) yields a new
algebraic finding: 7th roots of unity require an extension to
GF$(3^6)$ = GF$(729)$ \BM. The algebraic reason is compelling:
$x^7-1$ factors over GF$(3)$ as $(x-1)\cdot\phi_7$, where $\phi_7$
is \emph{irreducible of degree~6} --- a degree-6 extension is therefore
minimal and necessary.

GF$(729)^*$ has order $728 = 8\times7\times13$. The prime factors
yield three scales:
\begin{itemize}
  \item $13$: GF$(27)$ level --- quantum numbers, charges, generations \BM
  \item $7$: order-7 elements exist in G$(6)$ over GF$(9)$ (Doc.~341 §6,
    \texttt{pruef\_341\_root7\_gf9}) \BM; their reading as \emph{phases}
    $k\cdot360°/7$ is an $\mathbb{R}$-interpretation --- characteristic~3 has no
    angles, and $\cos(2\pi/7)\notin\mathrm{GF}(3^k)$ for any $k$ --- hence \SM
  \item $8$: octagonal structure; $360°/8 = 45°$ --- near
    $\theta_{23}^\text{PMNS}\approx49.2°$ (candidate \SM)
\end{itemize}

\textbf{Concrete observation \SM:}
$\theta_{23}^\text{PMNS}\approx49.2°$ lies $2.2°$ from $360°/7\approx51.4°$
--- within measurement uncertainty, but not a proof.
The CP phase $\delta_{CP}^\text{CKM}\approx68°$ has no direct relation to
7th roots of unity; $\delta_{CP}^\text{PMNS}\approx195°\approx4\times360°/7-10°$
is likewise not a direct match.

\textbf{What this means:} GF$(729)$ is the natural next step beyond GF$(27)$.
If CKM/PMNS angles reside at this level, Doc.~348 Theorem~D would need to be
executed on GF$(729)$ --- a concrete open task, not a structural obstacle.

\textbf{Why $P_+$ and not $P_-$ \SM:}
The chiral projection $P_+$ selects the left-handed Sd sector (A095).
Why nature selects $P_+$ (and not $P_-$) is geometrically plausible
but not derived from a deeper symmetry breaking \SM.

\section{Verification Script}

\begin{itemize}[leftmargin=2em,itemsep=2pt,topsep=2pt]
  \item Numerical verification of the SSB chain: $Q_\text{vac}=0$,
    residual symmetry $U(1)_Q$, $W^\pm/Z^0$ mass ratio from $\theta_W$,
    Goldstone counting, GF$(729)$ order:
    \texttt{python/Dok355\_Skripte/pruef\_355\_ssb.py}
\end{itemize}

\section{Use of AI Tools}

This document was produced with the assistance of AI-based tools.
The questions, postulates and all content decisions originate with the author;
the tools formulate, compute and generate verification scripts.
Responsibility for content and errors lies entirely with the author.
The detailed wording is in A010.
